% Madelung transform: shows equivalance of Schroedinger to the Continuity equation
%
%
\section{Madelung transform}
\label{app_mtrans}
This appendix provides a full derivation of the continuity and Bernoulli (Euler) equations from the Schr\"odinger equation via the Madelung transform. This is similar to the translated Madelung paper (Appendix \ref{app_m}) but uses notation consistent with this thesis and with Short \citep{bib_short}. The role of the pressure term is discussed elsewhere and is extracted out from this derivation in the appropriate place.
\begin{equation}
i\nu\frac{\partial\psi}{\partial t}= -\frac{\nu^2}{2}\nabla^2\psi + V\psi
\end{equation}

\noindent \textbf{Convention:} This derivation uses the positive-phase convention $\psi = \alpha\, e^{+i\varphi/\nu}$, so the velocity is $v = +\nabla\varphi$. In the main body of the thesis (Chapters \ref{ch_3}--\ref{ch_5}) the negative-phase convention $\psi = \alpha\, e^{-i\varphi_v/\nu}$ is used, giving $v = -\nabla\varphi_v$. The two are related by $\varphi = -\varphi_v$; all physical results are the same.
\begin{equation}
 \psi = (1+ \delta)^{1/2}exp\left(\frac{i\varphi}{\nu}\right) = \alpha exp\left(\frac{i\varphi}{\nu}\right)
\end{equation}

\noindent LHS, Differentiate $\psi$ with respect to time :
\begin{equation}
\frac{\partial\psi}{\partial t} = \frac{\partial\alpha}{\partial t}exp\left(\frac{i\varphi}{\nu}\right) + \frac{i\alpha}{\nu}exp\left(\frac{i\varphi}{\nu}\right)\frac{\partial\varphi}{\partial t}
\end{equation}

\begin{equation}
i\nu\frac{\partial\psi}{\partial t}=i\nu\frac{\partial\alpha}{\partial t}exp\left(\frac{i\varphi}{\nu}\right)-\psi\frac{\partial\varphi}{\partial t}
\end{equation}

\noindent Tidy up:
\begin{equation}
LHS = \frac{i\nu\psi}{\alpha}\frac{\partial\alpha}{\partial t} - \frac{\partial\varphi}{\partial t}
\end{equation}

\noindent Now for RHS, find $\nabla^2 \psi$
\begin{eqnarray}
\nabla\psi &=& \nabla\left(\alpha exp\left(\frac{i\varphi}{\nu}\right)\right)\\
&=& (\nabla\alpha) exp\left(\frac{i\varphi}{\nu}\right) + \frac{i\alpha}{\nu} exp\left(\frac{i\varphi}{\nu}\right)\nabla\varphi \nonumber
\end{eqnarray}

\begin{eqnarray}
\nabla^2\psi &=& \nabla . (\nabla \psi) \\ \nonumber
&=& \nabla.[(\nabla\alpha) exp\left(\frac{i\varphi}{\nu}\right) + \frac{i\alpha}{\nu} exp\left(\frac{i\varphi}{\nu}\right)\nabla\varphi] \\ \nonumber
&=& (\nabla^2\alpha)exp\left(\frac{i\varphi}{\nu}\right) +\frac{i}{\nu} exp\left(\frac{i\varphi}{\nu}\right)\nabla\varphi.\nabla\alpha \\ \nonumber
&+&\frac{i\alpha}{\nu} exp\left(\frac{i\varphi}{\nu}\right)\nabla^2\varphi +\frac{i}{\nu} \nonumber exp\left(\frac{i\varphi}{\nu}\right)\nabla\varphi.\nabla\alpha \\ \nonumber
&+& \frac{i^2\alpha}{\nu^2}exp\left(\frac{i\varphi}{\nu}\right)\nabla\varphi.\nabla\varphi \nonumber
\end{eqnarray}

\noindent Note the expressions for the potential term (Bernoulli equation, here $V$ denotes the external gravitational potential, not the effective potential of Ch4) and the pressure-like term $P$:
\begin{equation}
V = -\frac{\partial\varphi}{\partial t}\psi - \frac{1}{2}\psi(\nabla\varphi)^2 
\end{equation}

\begin{equation}
P = \frac{\nu^2}{2}\frac{(\nabla^2\alpha)}{\alpha}
\end{equation}

\noindent Leads to (RHS = K + V):
\begin{eqnarray}
RHS &=& -\frac{\nu^2\psi}{2\alpha}(\nabla^2\alpha) -\frac{i\nu\psi}{\alpha}(\nabla\varphi).(\nabla\alpha) \\ \nonumber
&-&\frac{i\nu}{2}\psi\nabla^2\varphi + \frac{1}{2}\psi(\nabla\varphi)^2 \\ \nonumber
&-&\frac{\partial\varphi}{\partial t}\psi - \frac{1}{2}\psi(\nabla\varphi)^2  \\ \nonumber
\end{eqnarray}

\noindent The last two terms above are from the potential. Now compare LHS and RHS, cancelling terms. The pressure and potential terms have been omitted here.
\begin{equation}
\frac{i\nu\psi}{\alpha}\frac{\partial\alpha}{\partial t} = -\frac{i\nu\psi}{\alpha}(\nabla\varphi).(\nabla\alpha)-\frac{i\nu}{2}\psi\nabla^2\varphi
\end{equation}

\noindent Which yields the continuity equation and hence showing that the Madelung transformation casts the Schr\"odinger equation into fluid dynamic form (Madelung: `hydrodynamischer').
\begin{equation}
\frac{\partial\alpha}{\partial t} +(\nabla\varphi).(\nabla\alpha)+\frac{\alpha}{2}\psi\nabla^2\varphi = 0
\end{equation}

\noindent To show that this is the continuity equation, work from the continuity equation until arriving at the above equation:
\begin{equation}
\frac{\partial\rho}{\partial t} +\nabla.(\rho\nabla\varphi) = 0
\end{equation}
\noindent Re-call the following definitions:
\begin{equation}
\rho = \psi\psi^* = |\psi|^2 = \alpha^2
\end{equation}

\noindent Then re-write the continuity equation as:
\begin{equation}
\frac{\partial\alpha^2}{\partial t} + \nabla . (\alpha^2\nabla\varphi)= 0
\end{equation}
\begin{equation}
2\alpha\frac{\partial\alpha}{\partial t} + \alpha^2\nabla^2\varphi + 2\alpha\nabla\alpha.\nabla\varphi = 0
\end{equation}
\noindent Finally leading to the desired form of the continuity equation that one arrives at when inserting the Madelung transformation into the Schr\"odinger equation:
\begin{equation}
\frac{\partial\alpha}{\partial t} + \frac{\alpha}{2}\nabla^2\varphi + \nabla\alpha.\nabla\varphi = 0
\end{equation}

\noindent This confirms that the Madelung transform casts the Schr\"odinger equation into the continuity equation (above) and the Bernoulli/Hamilton-Jacobi equation (from the real part, yielding the potential and pressure terms). Together these form the Madelung pair, which is the fluid-dynamical equivalent of the Schr\"odinger equation. The pressure-like term $P$ is discussed in detail in Chapter \ref{ch_3}.

